% Part: many-valued-logic
% Chapter: three-valued-logics
% Section: lukasiewicz

\documentclass[../../../include/open-logic-section]{subfiles}

\begin{document}

\olfileid{mvl}{inf}{god}

\olsection{گوڈل دیاں منطقاں}

\begin{editorial}
  ایہ لامتناہی قدراں والی گوڈل دی منطق بارے حصے دا نکا جیہا مسودہ اے۔
\end{editorial}

\begin{defn}\ollabel{def:goedel}
لامتناہی قدراں والی گوڈل دی منطق~$\LogGod[\infty]$ ایس میٹرکس
نال متعین اے۔
\begin{enumerate}
  \item معیاری قضیاتی زبان $\Lang L_0$ جیہدے وچ
  $\lfalse$، $\lnot$، $\land$، $\lor$، $\lif$ نیں۔
  \item صداقتی قدراں دا سیٹ $V_\infty$۔
  \item $1$ اکّو مقرر قدر اے، یعنی $V^+ = \{1\}$۔
  \item صداقتی فنکشن ہِٹھ دِتّے نیں۔
  \begin{align*}
    \tf{\lfalse} & = 0\\
    \tf{\lnot}[\LogGod](x) & = \begin{cases}
      $1$ & \text{جے } x =0\\
      $0$ & \text{ہور صورت وچ}
    \end{cases}\\
    \tf{\land}[\LogGod](x,y) & = \min(x,y)\\
    \tf{\lor}[\LogGod](x,y) & = \max(x,y)\\
    \tf{\lif}[\LogGod](x,y) & = \begin{cases}
      1 & \text{جے } x \le y\\
      y & \text{ہور صورت وچ۔}
    \end{cases}
    \end{align*}
\end{enumerate}
$m$ قدراں والی گوڈل دی منطق ایسے طرح متعین اے، بس $V = V_m$ اے۔
\end{defn}

\begin{prop}
  \olref[thr][god]{defn:goedel} وچ متعین منطق $\LogGod[3]$،
  \olref{def:goedel} وچ متعین منطق $\LogGod[3]$ دے عین برابر اے۔
\end{prop}

\begin{proof}
  ایہ گل \olref[thr][god]{defn:goedel} وچ دِتّے رابطیاں دے
  صداقتی جدولاں نوں \olref{def:goedel} دیاں مساواتاں توں ملن
  والے جدولاں نال ملا کے دِسدی اے۔
  \begin{center}
    \begin{tabular}{c|c} 
      $\tf{\lnot}[\LogGod[3]]$ & \\ 
      \hline  
      $1$ & $0$ \\ 
      $1/2$ & $0$ \\
      $0$ & $1$ 
    \end{tabular}
    \quad
    \begin{tabular}{c|ccc} 
      $\tf{\land}[\LogGod]$ & $1$ & $1/2$ & $0$ \\ 
      \hline 
      $1$ & $1$ & $1/2$ & $0$ \\ 
      $1/2$ & $1/2$ & $1/2$ & $0$\\ 
      $0$ & $0$ & $0$ & $0$ 
    \end{tabular}
    \\[2ex]
    \begin{tabular}{c|ccc} 
      $\tf{\lor}[\LogGod]$ & $1$ & $1/2$ & $0$ \\ 
      \hline 
      $1$ & $1$ & $1$ & $1$ \\ 
      $1/2$ & $1$ & $1/2$ & $1/2$ \\
      $0$ & $1$ & $1/2$ & $0$ 
    \end{tabular}
    \quad
    \begin{tabular}{c|ccc} 
      $\tf{\lif}[\LogGod]$ & $1$ & $1/2$ & $0$ \\ 
      \hline 
      $1$ & $1$ & $1/2$ & $0$ \\ 
      $1/2$ & $1$ & $1$ & $0$  \\ 
      $0$ & $1$ & $1$ & $1$ 
    \end{tabular}
  \end{center}
\end{proof}

\begin{prop}\ollabel{prop:god-infty-m}
  جے $\Gamma \Entails[\LogGod[\infty]] !B$، تاں $\Gamma
  \Entails[\LogGod[m]] !B$ ہر~$m \ge 2$ لئی۔
\end{prop}

\begin{proof}
  مشق۔
\end{proof}

\begin{prob}
  \olref[mvl][inf][god]{prop:god-infty-m} ثابت کرو۔
\end{prob}

اصل وچ اُلٹی گل وی قائم اے۔

$\LogGod[3]$ وانگ، $\LogGod[\infty]$ وچ اوہ سارے !!{formula}s
ہمیشہ صادق نیں جیہڑے وجدانیت پسند منطق وچ معتبر نیں۔ اوہی
مثالاں جیہڑیاں ہمیشہ صادق نہیں سن، ~$\LogGod[\infty]$ وچ وی
ہمیشہ صادق نہیں۔
\begin{align*}
  & p \lor \lnot p && (p \lif q) \lif (\lnot p \lor q) \\
  & \lnot\lnot p \lif p && \lnot(\lnot p \land \lnot q) \lif (p \lor q) \\
  & ((p \lif q) \lif p) \lif p && \lnot(p \lif q) \lif (p \land \lnot q)
\end{align*}
اوہ !!{formula} جیہڑا وجدانیت پسند منطق وچ معتبر نہیں پر
~$\LogGod[3]$ وچ ہمیشہ صادق اے، یعنی $(p \lif q) \lor (q \lif
p)$، ~$\LogGod[\infty]$ وچ وی ہمیشہ صادق اے۔ دراصل
$\LogGod[\infty]$ نوں وجدانیت پسند منطق وچ سانچا
$(!A \lif !B) \lor (!B \lif !A)$ شامل کر کے حاصل کر سکدے آں۔
ایہ گل مائیکل ڈمٹ نے دکھائی سی۔ ایس لئی $\LogGod[\infty]$
نوں اکثر گوڈل–ڈمٹ منطق~$\Log{LC}$ آکھدے آں۔

\begin{prob}
  دکھاؤ کہ $(p \lif q) \lor (q \lif p)$، ~$\LogGod[\infty]$
  وچ ہمیشہ صادق اے۔
\end{prob}

\begin{prob}
  دکھاؤ کہ $(p \lif q) \lor (q \lif r) \lor (r \lif s)$،
  جیہڑا $\LogGod[3]$ وچ ہمیشہ صادق اے، ~$\LogGod[\infty]$ وچ
  ہمیشہ صادق نہیں۔
\end{prob}

\end{document}
